% Chinese questions and concise solutions. Compile using XeLaTeX.
\documentclass[UTF8,10pt,letterpaper,fontset=fandol]{ctexart}
\usepackage[margin=1in]{geometry}
\usepackage{amsmath,amssymb}
\usepackage{needspace}
\usepackage[unicode,hidelinks]{hyperref}
\pagestyle{plain}
\raggedbottom
\setlength{\parindent}{0pt}
\setlength{\parskip}{1pt}
\setlength{\emergencystretch}{2em}
\setcounter{secnumdepth}{0}
\ctexset{section={format=\Large\bfseries,beforeskip=1.2ex,afterskip=0.8ex}}
\newcommand{\F}{\mathbb F}
\newcommand{\Solution}{\par\Needspace{4\baselineskip}\textbf{解答}\par}
\begin{document}
\begin{center}
{\LARGE 作业 4}\\[0.6em]
2026 年 10 月 16 日\\[0.25em]
姓名：YuanHao Tang\quad 学号：3120265621
\end{center}

\section{第 1 题}
求下列 $\lambda$-矩阵的 Smith 标准形：
\[
A(\lambda)=\begin{pmatrix}
1-\lambda&\lambda^2&\lambda\\
2\lambda-1&\lambda(1-\lambda)&-2\lambda\\
1+\lambda^2&\lambda^2&-\lambda^2
\end{pmatrix}.
\]
\Solution
记 $D_k$ 为全部 $k$ 阶子式的首一最大公因式，$D_0=1$，则不变多项式为 $d_k=D_k/D_{k-1}$。
由 $2(1-\lambda)+(2\lambda-1)=1$，得 $D_1=1$。第二、三列都含因子 $\lambda$，故所有二阶子式均被 $\lambda$ 整除，而
\[
\begin{vmatrix}1-\lambda&\lambda\\2\lambda-1&-2\lambda\end{vmatrix}=-\lambda,
\]
所以 $D_2=\lambda$。又 $\det A=-\lambda^2(\lambda+1)$，故 $D_3=\lambda^2(\lambda+1)$。因此
\[
\boxed{S_A(\lambda)=\operatorname{diag}\bigl(1,\lambda,\lambda(\lambda+1)\bigr)}.
\]

\section{第 2 题}
设 $A(\lambda)$ 为下列 $n\times n$ 的 $\lambda$-矩阵：
\[
A(\lambda)=\begin{pmatrix}
\lambda-a&b_1&0&\cdots&0\\
0&\lambda-a&b_2&\ddots&\vdots\\
\vdots&\ddots&\ddots&\ddots&0\\
0&\cdots&0&\lambda-a&b_{n-1}\\
0&\cdots&\cdots&0&\lambda-a
\end{pmatrix},
\]
其中 $b_1,b_2,\ldots,b_{n-1}$ 均为非零常数。求 $A(\lambda)$ 的不变多项式。
\Solution
对 $1\le k<n$，取第 $1,\ldots,k$ 行和第 $2,\ldots,k+1$ 列，所得子式为非零常数
$b_1\cdots b_k$，故 $D_k=1$。又 $D_n=\det A=(\lambda-a)^n$，因此
\[
\boxed{d_1=\cdots=d_{n-1}=1,\qquad d_n=(\lambda-a)^n}.
\]

\newpage
\section{第 3 题}
求下列 $\lambda$-矩阵的 Smith 标准形：
\[
A(\lambda)=\begin{pmatrix}
\lambda&0&\cdots&0&a_0\\
-1&\lambda&\ddots&\vdots&a_1\\
0&-1&\ddots&0&\vdots\\
\vdots&\ddots&\ddots&\lambda&a_{n-2}\\
0&\cdots&0&-1&\lambda+a_{n-1}
\end{pmatrix}.
\]
\Solution
对 $1\le k<n$，第 $2,\ldots,k+1$ 行与第 $1,\ldots,k$ 列构成的子式为 $(-1)^k$，故 $D_k=1$。
按最后一列展开可得
\[
D_n=\det A=\lambda^n+a_{n-1}\lambda^{n-1}+\cdots+a_1\lambda+a_0=:p(\lambda).
\]
因此
\[
\boxed{S_A(\lambda)=\operatorname{diag}\bigl(\underbrace{1,\ldots,1}_{n-1},p(\lambda)\bigr)}.
\]

\section{第 4 题}
对矩阵 $A\in\F^{n\times n}$ 及不定元 $\lambda$，称 $\lambda E-A$ 为 $A$ 的特征矩阵，其中 $E$ 是单位矩阵。求下列矩阵的特征矩阵的 Smith 标准形：
\[
A=\begin{pmatrix}-7&-12&-6\\3&5&3\\3&6&2\end{pmatrix}.
\]
\Solution
当域的特征不为 $3$（包括实数域和复数域）时，令
\[
B(\lambda)=\lambda I-A=\begin{pmatrix}
\lambda+7&12&6\\-3&\lambda-5&-3\\-3&-6&\lambda-2
\end{pmatrix}.
\]
因含非零常数 $-3$，$D_1=1$。$B(-1)$ 的秩为 $1$，故全部二阶子式含因子 $\lambda+1$，而
\[
\begin{vmatrix}\lambda+7&6\\-3&-3\end{vmatrix}=-3(\lambda+1),
\]
所以 $D_2=\lambda+1$。又 $D_3=\det B=(\lambda+1)^2(\lambda-2)$，故
\[
\boxed{S_B(\lambda)=\operatorname{diag}\bigl(1,\lambda+1,(\lambda+1)(\lambda-2)\bigr)}.
\]
若域的特征为 $3$，则 $A=-I$，标准形为 $\boxed{(\lambda+1)I_3}$。
\end{document}
